\chapter{Introduction}
\label{chap:introduction}

\section{Background}
\label{section:background}

Figure skating captivates millions through Winter Olympics and ISU Grand Prix events, blending athleticism with artistry. The sport features six primary jumps—toe loop, flip, Lutz (toe-assisted takeoff) and Salchow, loop, Axel (edge takeoff)—distinguished by subtle takeoff mechanics that challenge even experienced viewers. During live competitions, technical panels use instant replays for unclear calls while broadcasters rely on pre-written program sheets that often misidentify jumps, leaving casual audiences confused about scoring decisions.

\section{Problem Statement}
\label{section:problem-statement}

Figure skating jumps often look very similar to the casual viewer, yet they receive different technical scores. This creates a major challenge for the audience because broadcasts often lack clear commentary or provide incorrect information. Even professionals find it difficult to judge edge usage and rotation in real time. Unlike objective sports like tennis or badminton, where a ball is clearly in or out, skating scores depend on how judges interpret each move. While technical panels use replays to fix errors, viewers are often left with misidentified jumps from pre-written program sheets. This confusion makes it hard for new fans to understand the scoring and reduces their overall interest in the sport.

\section{Solution Overview}
\label{section:solution-overview}

Figure Flow processes broadcast video through a computer vision pipeline to deliver real-time jump identification. The system detects skaters with YOLO, extracts pose keypoints using MediaPipe Pose, and classifies jumps with LSTM temporal analysis. Phase 1 handles Axel vs Non-Axel classification; Phase 2 extends to Toe vs Edge jumps.

\subsection{Features}
\label{subsection:features}

\begin{enumerate}[leftmargin=80pt]
    \item Skater Detection: YOLO tracks skater across video frames
    \item Pose Extraction: MediaPipe Pose captures leg position, body alignment, rotation
    \item Axel Classification: LSTM model distinguishes forward takeoff Axel jumps (Phase 1)
    \item Toe vs Edge: Extended classification for toe pick vs pure edge takeoffs (Phase 2)
    \item Visualization Interface: On-screen jump labels with simple explanations
\end{enumerate}

\section{Target User}
\label{section:target-user}

Figure Flow serves three primary user groups:\\

\textbf{1. Casual Viewers (Primary)}: Sports fans seeking accessible explanations during Olympics/Grand Prix broadcasts. Minimal skating knowledge; frustrated by technical jargon.

\textbf{2. Serious Enthusiasts}: Dedicated fans who follow ISU circuits and want real-time verification of judging calls. Technical interest but lack expert-level jump recognition skills.

\textbf{3. Non-Expert Commentary}: Commentary teams and online platforms (Twitch, YouTube) needing reliable jump identification beyond program sheets. Requires low-latency integration with existing video pipelines.

\section{Terminology}
\label{section:terminology}

\begin{description}
    \item[Toe Jump] Uses toe pick for takeoff assistance (toe loop, flip, Lutz)
    \item[Edge Jump] Uses skating edge only for takeoff (Salchow, loop, Axel)  
    \item[Axel Jump] Unique forward takeoff edge jump (2.5 or 3.5 rotations)
    \item[Pose Keypoints] Body joint coordinates (hips, knees, ankles) from MediaPipe
    \item[Temporal Sequence] Time-series pose data across multiple frames
    \item[Technical Panel] ISU officials identifying jump types live
\end{description}